\documentclass[11pt,a4paper]{article}

% Packages
\usepackage[utf8]{inputenc}
\usepackage[T1]{fontenc}
\usepackage[margin=2.5cm]{geometry}
\usepackage{graphicx}
\usepackage{xcolor}
\usepackage{listings}
\usepackage{hyperref}
\usepackage{fancyhdr}
\usepackage{tcolorbox}
\usepackage{enumitem}
\usepackage{tikz}

% TikZ libraries
\usetikzlibrary{shapes,arrows,positioning}

% Code listing style
\lstset{
    language=C,
    basicstyle=\ttfamily\small,
    keywordstyle=\color{blue}\bfseries,
    commentstyle=\color{green!60!black}\itshape,
    stringstyle=\color{red},
    showstringspaces=false,
    breaklines=true,
    frame=single,
    numbers=left,
    numberstyle=\tiny\color{gray},
    tabsize=2
}

% Color boxes
\tcbuselibrary{skins,breakable}
\newtcolorbox{objectivebox}{
    colback=blue!5!white,
    colframe=blue!75!black,
    title=Lab Objectives,
    fonttitle=\bfseries
}

\newtcolorbox{prereqbox}{
    colback=yellow!5!white,
    colframe=yellow!75!black,
    title=Prerequisites,
    fonttitle=\bfseries
}

\newtcolorbox{solutionbox}{
    colback=green!5!white,
    colframe=green!75!black,
    title=Solution,
    fonttitle=\bfseries,
    breakable
}

% Header and footer
\pagestyle{fancy}
\fancyhf{}
\lhead{USP Training Course 2}
\rhead{Lab 2.1: Event-Driven Programming}
\cfoot{\thepage}

\title{
    \textbf{Lab 2.1: Event-Driven Programming with LBM}\\
    \large Course 2: LoRa Basics Modem Architecture
}
\author{USP Training Series}
\date{}

\begin{document}

\mailtitle
\thispagestyle{fancy}

% ============================================================================
\section{Lab Overview}
% ============================================================================

\begin{objectivebox}
By the end of this lab, you will be able to:
\begin{itemize}
    \item Understand the LBM event-driven architecture
    \item Implement event handlers for LoRaWAN events
    \item Handle JOINED, TXDONE, DOWNDATA events
    \item Process event queues correctly
    \item Track application state based on modem events
\end{itemize}
\end{objectivebox}

\vspace{0.5cm}

\begin{prereqbox}
\textbf{Required Knowledge:}
\begin{itemize}
    \item Completion of Lab 1.1 (Basic LoRaWAN Join)
    \item Understanding of event-driven programming
    \item Familiarity with LBM API
\end{itemize}

\vspace{0.3cm}

\textbf{Hardware Required:}
\begin{itemize}
    \item Xiao nRF54L15 + LR1120 board
    \item USB cable
    \item LoRaWAN gateway in range
\end{itemize}
\end{prereqbox}

% ============================================================================
\section{Lab Duration}
% ============================================================================

\textbf{Estimated Time:} 60 minutes

% ============================================================================
\section{Background}
% ============================================================================

\subsection{LBM Event-Driven Architecture}

LoRa Basics Modem (LBM) uses an event-driven architecture where the application responds to asynchronous events from the modem. This design:

\begin{itemize}
    \item Decouples application logic from modem timing
    \item Enables low-power operation (sleep between events)
    \item Simplifies handling of asynchronous LoRaWAN operations
    \item Supports multiple concurrent operations
\end{itemize}

\subsection{Event Types}

Common LBM events include:

\begin{center}
\begin{tabular}{|l|p{8cm}|}
\hline
\textbf{Event} & \textbf{Description} \\
\hline
JOINED & Device successfully joined network \\
TXDONE & Uplink transmission completed \\
DOWNDATA & Downlink data received \\
UPLOAD\_DONE & File upload completed (FUOTA) \\
LINK\_CHECK & Link check result available \\
ALARM & Scheduled alarm triggered \\
\hline
\end{tabular}
\end{center}

% ============================================================================
\section{Lab Instructions}
% ============================================================================

\subsection{Step 1: Initialize LBM Stack}

Create a new application that uses LBM instead of Zephyr's native LoRaWAN stack.

\subsection{Step 2: Implement Event Handler}

Create an event handler function that processes all modem events in a loop.

\subsection{Step 3: Handle JOIN Event}

When JOINED event is received, initiate first uplink transmission.

\subsection{Step 4: Handle TXDONE Event}

Track transmission status (confirmed/unconfirmed, ACK received, etc.).

\subsection{Step 5: Handle DOWNDATA Event}

Process received downlink messages and extract payload.

\subsection{Step 6: Build and Test}

\begin{lstlisting}[language=bash]
west build -b xiao_nrf54l15 .
west flash
\end{lstlisting}

% ============================================================================
\section{Solution}
% ============================================================================

\begin{solutionbox}

\subsection{Complete Implementation}

\begin{lstlisting}
#include <zephyr/kernel.h>
#include <zephyr/logging/log.h>
#include <smtc_modem_api.h>
#include <smtc_board.h>

LOG_MODULE_REGISTER(lab_2_1, LOG_LEVEL_INF);

/* LoRaWAN Credentials */
#define LORAWAN_DEV_EUI    { 0x00, 0x00, 0x00, 0x00, \
                             0x00, 0x00, 0x00, 0x01 }
#define LORAWAN_JOIN_EUI   { 0x00, 0x00, 0x00, 0x00, \
                             0x00, 0x00, 0x00, 0x00 }
#define LORAWAN_APP_KEY    { 0x00, 0x00, 0x00, 0x00, \
                             0x00, 0x00, 0x00, 0x00, \
                             0x00, 0x00, 0x00, 0x00, \
                             0x00, 0x00, 0x00, 0x00 }

/* Application state */
static struct {
    bool joined;
    uint32_t uplink_count;
    uint32_t downlink_count;
    int16_t last_rssi;
    int8_t last_snr;
} app_state = {0};

/**
 * @brief Handle JOINED event
 */
static void handle_joined_event(void)
{
    LOG_INF("===========================================");
    LOG_INF("JOINED EVENT: Network joined successfully");
    LOG_INF("===========================================");

    app_state.joined = true;

    /* Get and display DevAddr */
    uint32_t dev_addr;
    smtc_modem_get_dev_addr(0, &dev_addr);
    LOG_INF("DevAddr: 0x%08X", dev_addr);

    /* Send first uplink */
    uint8_t payload[] = "Hello from LBM!";
    smtc_modem_request_uplink(0, 2, false,
                              payload, sizeof(payload));
    LOG_INF("First uplink queued");
}

/**
 * @brief Handle TXDONE event
 */
static void handle_txdone_event(smtc_modem_event_t *event)
{
    app_state.uplink_count++;

    LOG_INF("TXDONE EVENT #%u:", app_state.uplink_count);
    LOG_INF("  Status: %s",
            event->event_data.txdone.status ==
            SMTC_MODEM_EVENT_TXDONE_CONFIRMED
                ? "CONFIRMED" : "UNCONFIRMED");

    if (event->event_data.txdone.status ==
        SMTC_MODEM_EVENT_TXDONE_CONFIRMED) {
        LOG_INF("  ACK received");
    }

    /* Schedule next uplink */
    k_work_schedule(&uplink_work, K_SECONDS(60));
}

/**
 * @brief Handle DOWNDATA event
 */
static void handle_downdata_event(smtc_modem_event_t *event)
{
    uint8_t rx_payload[255];
    uint8_t rx_payload_size;
    smtc_modem_dl_metadata_t rx_metadata;

    app_state.downlink_count++;

    LOG_INF("DOWNDATA EVENT #%u:", app_state.downlink_count);

    /* Get downlink data */
    smtc_modem_get_downlink_data(
        0, rx_payload, &rx_payload_size,
        &rx_metadata);

    LOG_INF("  Port: %d", rx_metadata.fport);
    LOG_INF("  RSSI: %d dBm", rx_metadata.rssi);
    LOG_INF("  SNR: %d dB", rx_metadata.snr);
    LOG_INF("  Length: %d bytes", rx_payload_size);

    if (rx_payload_size > 0) {
        LOG_HEXDUMP_INF(rx_payload, rx_payload_size,
                       "Payload");

        /* Store link quality */
        app_state.last_rssi = rx_metadata.rssi;
        app_state.last_snr = rx_metadata.snr;
    }
}

/**
 * @brief Main event processing loop
 */
static void process_modem_events(void)
{
    smtc_modem_event_t event;
    uint8_t pending_events;

    /* Process all pending events */
    while (smtc_modem_get_event(&event, &pending_events)
           == SMTC_MODEM_RC_OK) {

        LOG_DBG("Event received: type=%d, pending=%d",
                event.event_type, pending_events);

        switch (event.event_type) {
        case SMTC_MODEM_EVENT_RESET:
            LOG_INF("RESET EVENT: Modem reset");
            break;

        case SMTC_MODEM_EVENT_ALARM:
            LOG_INF("ALARM EVENT: Timer expired");
            break;

        case SMTC_MODEM_EVENT_JOINED:
            handle_joined_event();
            break;

        case SMTC_MODEM_EVENT_TXDONE:
            handle_txdone_event(&event);
            break;

        case SMTC_MODEM_EVENT_DOWNDATA:
            handle_downdata_event(&event);
            break;

        case SMTC_MODEM_EVENT_JOINFAIL:
            LOG_ERR("JOIN FAIL EVENT");
            /* Retry join */
            k_work_schedule(&join_work, K_SECONDS(10));
            break;

        case SMTC_MODEM_EVENT_LINK_CHECK:
            LOG_INF("LINK CHECK EVENT:");
            LOG_INF("  Margin: %d dB",
                    event.event_data.link_check.margin);
            LOG_INF("  Gateways: %d",
                    event.event_data.link_check.gw_cnt);
            break;

        case SMTC_MODEM_EVENT_MUTE:
            LOG_WRN("MUTE EVENT: Device muted by network");
            break;

        default:
            LOG_WRN("Unhandled event: %d",
                    event.event_type);
            break;
        }
    }
}

/* Work queue for periodic uplinks */
static void uplink_work_handler(struct k_work *work)
{
    if (!app_state.joined) {
        LOG_WRN("Not joined, skipping uplink");
        return;
    }

    uint8_t payload[32];
    int len = snprintf(payload, sizeof(payload),
                      "Uplink #%u", app_state.uplink_count + 1);

    smtc_modem_request_uplink(0, 2, false, payload, len);
    LOG_INF("Uplink queued: %s", payload);
}

K_WORK_DELAYABLE_DEFINE(uplink_work, uplink_work_handler);

/* Work queue for join retry */
static void join_work_handler(struct k_work *work)
{
    LOG_INF("Retrying join...");
    smtc_modem_join_network(0);
}

K_WORK_DELAYABLE_DEFINE(join_work, join_work_handler);

int main(void)
{
    int ret;

    LOG_INF("==============================================");
    LOG_INF("Lab 2.1: Event-Driven Programming with LBM");
    LOG_INF("==============================================");

    /* Initialize board-specific hardware */
    smtc_board_init();

    /* Initialize modem */
    ret = smtc_modem_init();
    if (ret != SMTC_MODEM_RC_OK) {
        LOG_ERR("Modem init failed: %d", ret);
        return -1;
    }

    LOG_INF("Modem initialized");

    /* Set credentials */
    uint8_t dev_eui[] = LORAWAN_DEV_EUI;
    uint8_t join_eui[] = LORAWAN_JOIN_EUI;
    uint8_t app_key[] = LORAWAN_APP_KEY;

    smtc_modem_set_deveui(0, dev_eui);
    smtc_modem_set_joineui(0, join_eui);
    smtc_modem_set_nwkkey(0, app_key);

    /* Set region */
    smtc_modem_set_region(0, SMTC_MODEM_REGION_EU_868);

    LOG_INF("Credentials configured");

    /* Start join procedure */
    LOG_INF("Starting OTAA join...");
    ret = smtc_modem_join_network(0);
    if (ret != SMTC_MODEM_RC_OK) {
        LOG_ERR("Join failed: %d", ret);
        return -1;
    }

    /* Main event loop */
    LOG_INF("Entering event loop");

    while (1) {
        /* Process modem events */
        process_modem_events();

        /* Run modem engine */
        smtc_modem_run_engine();

        /* Sleep until next event */
        k_sleep(K_MSEC(100));
    }

    return 0;
}
\end{lstlisting}

\subsection{CMakeLists.txt}

\begin{lstlisting}[language=bash]
cmake_minimum_required(VERSION 3.20.0)
find_package(Zephyr REQUIRED HINTS $ENV{ZEPHYR_BASE})
project(lab_2_1)

target_sources(app PRIVATE src/main.c)

# Include LBM library
target_link_libraries(app PRIVATE lora_basics_modem)
\end{lstlisting}

\subsection{prj.conf}

\begin{lstlisting}[language=bash]
# LBM Configuration
CONFIG_LORA_BASICS_MODEM=y
CONFIG_LBM_REGION_EU868=y

# Logging
CONFIG_LOG=y
CONFIG_LBM_LOG_LEVEL_DBG=y
CONFIG_LOG_BUFFER_SIZE=4096

# Stack sizes
CONFIG_MAIN_STACK_SIZE=4096
CONFIG_SYSTEM_WORKQUEUE_STACK_SIZE=2048
\end{lstlisting}

\end{solutionbox}

% ============================================================================
\section{Expected Output}
% ============================================================================

\begin{lstlisting}
==============================================
Lab 2.1: Event-Driven Programming with LBM
==============================================
[00:00:00.100] <inf> lab_2_1: Modem initialized
[00:00:00.105] <inf> lab_2_1: Credentials configured
[00:00:00.110] <inf> lab_2_1: Starting OTAA join...
[00:00:00.115] <inf> lab_2_1: Entering event loop
[00:00:05.234] <inf> lab_2_1: ===================
[00:00:05.235] <inf> lab_2_1: JOINED EVENT
[00:00:05.236] <inf> lab_2_1: ===================
[00:00:05.237] <inf> lab_2_1: DevAddr: 0x260B1234
[00:00:05.238] <inf> lab_2_1: First uplink queued
[00:00:06.456] <inf> lab_2_1: TXDONE EVENT #1:
[00:00:06.457] <inf> lab_2_1:   Status: UNCONFIRMED
[00:01:06.457] <inf> lab_2_1: Uplink queued: Uplink #2
[00:01:07.678] <inf> lab_2_1: TXDONE EVENT #2:
[00:01:07.679] <inf> lab_2_1:   Status: UNCONFIRMED
[00:01:08.123] <inf> lab_2_1: DOWNDATA EVENT #1:
[00:01:08.124] <inf> lab_2_1:   Port: 2
[00:01:08.125] <inf> lab_2_1:   RSSI: -45 dBm
[00:01:08.126] <inf> lab_2_1:   SNR: 10 dB
[00:01:08.127] <inf> lab_2_1:   Length: 4 bytes
\end{lstlisting}

% ============================================================================
\section{Deliverables}
% ============================================================================

\begin{enumerate}[itemsep=10pt]
    \item \textbf{Source Code} with event handlers for all required events
    \item \textbf{Serial Log} showing successful event processing
    \item \textbf{Event Timing Analysis}
    \begin{itemize}
        \item Time from join to first TXDONE
        \item Average time between uplinks
        \item Response time to downlinks
    \end{itemize}
    \item \textbf{Lab Report} describing event flow and state transitions
\end{enumerate}

% ============================================================================
\section{Additional Challenges}
% ============================================================================

\subsection{Challenge 1: Event Statistics}

Track statistics for each event type:
\begin{itemize}
    \item Count of each event received
    \item Average time between events
    \item Display summary every 10 minutes
\end{itemize}

\subsection{Challenge 2: Custom Event Actions}

Implement custom actions based on link quality:
\begin{itemize}
    \item If RSSI < -100 dBm, increase SF
    \item If SNR < 0 dB, request link check
    \item Log quality trends
\end{itemize}

\vfill

\noindent\rule{\textwidth}{0.4pt}
\begin{center}
\textit{USP Training Course 2 - Lab 2.1}\\
\textit{Version 1.0 - 2025}
\end{center}

\end{document}
